\section{资源约束、采样与显存优化}

\subsection{2 GiB 硬门禁}

每个训练或特征提取配置同时监测 allocator 已分配和预留峰值：
\begin{equation}
  M_{\max}=\max_t M_{\mathrm{a}}(t),\qquad
  R_{\max}=\max_t M_{\mathrm{r}}(t),
  \label{eq:memory}
\end{equation}
\begin{equation}
  \mathrm{feasible}
  =\mathbf{1}\left[
  M_{\max}\leq2048\ \mathrm{MiB}
  \land
  R_{\max}\leq2048\ \mathrm{MiB}
  \right].
  \label{eq:feasible}
\end{equation}
任何峰值超过预算的进程立即终止，不保留其精度结果作为可行配置。

\subsection{三种帧采样}

均匀采样把视频分为 $K$ 段，在每段中心附近取样，测试时不加抖动：
\begin{equation}
  t_k=\left\lfloor
  \left(k-\frac{1}{2}\right)\frac{T}{K}
  \right\rfloor.
  \label{eq:uniform}
\end{equation}
稠密采样在长度为 $(K-1)d+1$ 的连续窗口中每隔 $d$ 帧取一帧。它能保留短时动作，但可能遗漏
长视频的另一动作阶段。自适应采样先均匀解码 32 个候选帧，计算灰度帧差运动强度
\begin{equation}
  m_t=\frac{1}{HW}
  \left\lVert I_{t+1}-I_t\right\rVert_1,
  \label{eq:motion}
\end{equation}
再用“运动强度 + 时间多样性”的贪心目标选择 $K$ 帧。候选帧数固定为 32，额外解码和评分
时间分别记录为 $T_d$ 与 $T_s$，并与普通采样统一计入端到端吞吐。

\subsection{混合精度、检查点与梯度累积}

混合精度在前向和反向中使用 FP16 计算并保留 FP32 主权重与损失缩放\cite{amp}，以降低激活和中间张量
显存。梯度检查点只保存模块边界，在反向传播时重算中间激活\cite{checkpoint}，以计算换显存。资源不足时以
小批量 $B_s$ 和累积步数 $G$ 构造有效批量
\begin{equation}
  B_{\mathrm{eff}}=B_sG.
  \label{eq:accumulation}
\end{equation}
所有技术消融均在同一输入、同一模型和同一步优化流程下比较，避免把批量变化与精度变化混杂。

\subsection{时间与空间预算}

为隔离“时间更长”与“空间更大”的选择，时间优先和空间优先配置保持每视频帧像素成本相同：
\begin{equation}
  K_tH_tW_t=K_sH_sW_s.
  \label{eq:equal_cost}
\end{equation}
时间优先使用 $16\times64\times64$，空间优先使用 $4\times128\times128$，二者均为 65536
个空间像素帧，批大小同为 8。均衡配置使用主实验的 $8\times112\times112$。最终同时比较
准确率、峰值显存和吞吐，而不是只看一个指标。
